\documentclass{article}
\usepackage{amsmath,amssymb}
\usepackage[margin=1in]{geometry}
\allowdisplaybreaks
\usepackage[UTF8]{ctex}
\begin{document}

\[
\underset{(p_{i1}^{c},p_{i1}^{u})}{\operatorname{maximize}}
\quad \text{sum of expected profits in both periods},
\qquad i=1,2,\quad j\ne i.
\]
目标取自 Garcia、Janssen 和 Shopova 的论文 \emph{容量不确定下的动态定价}（2022 年 1 月 26 日）中未编号的完美贝叶斯均衡定义的文字部分，见 PDF 实际第 ~7 至 9 页。 每家公司在第一阶段的策略，在其产能类型给定的情况下是最优的，前提是考虑到竞争对手的策略$(p_{j1}^{c},p_{j1}^{u})$以及均衡状态下的后续反应。 策略可以混合使用。

遵循以下平衡要求：
\begin{enumerate}
\item 第 2 期开始时，未售出商品的企业 $i$ 在可能的情况下用贝叶斯法则更新其关于竞争对手容量受限的信念 $\theta(p_{j1})$。
\item 对于每一对第一时期的定价$(p_{11},p_{21})$，每个拥有未售出产能的厂商$i$都会选择$p_{i2}$来最大化第二时期的利润，这一决策是考虑到竞争对手在第二时期的定价$p_{j2}$，以及更新后的信念$\theta(p_{j1})$而做出的。
\item 每家公司$i$都选择$(p_{i1}^{c},p_{i1}^{u})$，以在给定$(p_{j1}^{c},p_{j1}^{u})$和均衡反应的情况下，最大化两个时期的预期利润。
\[
p_{i2}(p_{i1},p_{j1}),\qquad i=1,2,\quad j\ne i.
\]
\end{enumerate}
源代码的策略与信念对象为：
\[
\bigl(p_{i1}^{c},p_{i1}^{u},p_{i2}(p_{i1},p_{j1})\bigr),
\qquad i=1,2,\quad j\ne i,
\qquad \theta(p_{j1}).
\]
这里，$c$和$u$分别表示受限产能类型与非受限产能类型。非销售方$i$对公司$j$在观察到$j$第一期的价格后属于受限产能类型这一情况持有这样的信念$\theta(p_{j1})$。在仍有产能的情况下，公司的第二期定价并不取决于其初始产能是受限还是非受限（脚注14）。

固定的基线输入和时序如下。共有两家公司，并且有两个销售周期。
\[
i\in\{1,2\},\qquad t\in\{1,2\}.
\]
一个受限制的公司只能销售一个单位的产品；而一个不受限制的公司则可以在两个时期都满足需求。初始产能是私有的信息。固定的先决条件得到满足。
\[
\Pr(\text{firm }i\text{ is constrained})=\alpha,
\qquad
\Pr(\text{firm }i\text{ is unconstrained})=1-\alpha,
\qquad i=1,2,
\]
各公司的产能类型是相互独立的。生产成本被标准化为零。消费者具有短期行为，且需求是单位化的；每个时期的市场规模以及支付意愿都被标准化为$1$。每个时期都会有一个消费者出现，他们会观察当前的价格，然后从价格最低的商家购买商品，前提是该价格不超过$1$。否则，消费者就会离开，不再回来。如果两家商家都收取相同的价格，那么所有的需求都会流向同一家商家，而每家商家销售商品的概率都是$1/2$。

在周期1之前，各企业会先评估自身的产能。之后，各企业会选择$p_{i1}$。在首次采购之后，各企业会关注第一周期的价格，并了解哪家企业进行了销售，如果有的话。他们会更新自己的信念；只有那些没有出售出产的企业的才会选择$p_{i2}$。

\emph{打印采购门槛冲突。}物理设备上的操作设置。~8表示当最低价格“不大于1”时，就会发生采购。而该页面上的图1则显示
\[
p_1<1,\qquad p_2<1.
\]
这里的弱购买阈值保持不变，这与论文后续提到的销售价格和价格区间$1$是一致的；不会同时施加严格的阈值限制。基准定义并未规定全球性的价格区间。特别是，论文明确考虑了第一阶段的偏差$p_1>1$，因此，低于均衡支撑水平的价格并不是可行的价格范围。

以下连续公式适用于基线均衡构建（物理部分~12--13）。用$S$表示在第1期进行销售的厂商，用$N$表示未进行销售的厂商。用$F^{S}$和$F^{N}$分别表示它们在第二期的价格分布函数，其中$F^{S}$描述了在仍具有生产能力的情况下销售厂商的定价策略。
\[
\theta\in(0,1),
\]
公司$N$通常被认为具有生产能力，而公司$S$的库存可能已经售罄，概率为$\theta$。因此，公司$N$是一个垄断者，概率为$\theta$，同时也会面临竞争，概率为$1-\theta$。其预期利润表达式如下：
\[
\pi_2^{N}(p_2)
 =\theta p_2+(1-\theta)\bigl(1-F^{S}(p_2)\bigr)p_2,
\qquad p_2\leq 1,
\]
带有$\pi_2^{N}(1)=\theta$的数学符号，其顺序保持不变。所有数字都需保留原样。
\begin{equation}
F^{S}(p_2)=1-\frac{\theta(1-p_2)}{(1-\theta)p_2},
\qquad p_2\in[\theta,1).
\tag{1}
\end{equation}
在该价格区间内，不受限制的销售公司预计在第二阶段的利润为
\[
\bigl(1-F^{N}(p_2)\bigr)p_2.
\]
它的冷漠需要……
\begin{equation}
F^{N}(p_2)=1-\frac{\theta}{p_2},
\qquad p_2\in[\theta,1),
\tag{2}
\end{equation}
在$p_2=1$位置有一个大小为$\theta$的质量点，该质量点属于$N$公司。只要两家公司仍有能力运作，它们都能获得预期的第二时期利润$\theta$。在后续的计算中，$\theta$可以取决于销售公司在第一时期的价格，此时将其写作$\theta(p)$。

限制条件$0<\theta<1$属于这些混合延续公式。这些公式在$\theta=0$或$\theta=1$处不会被计算。论文中详细描述了完全揭示历史数据如何产生$0$或$1$这样的第二时期价格：当销售公司已知会保留产能时，就会出现竞争；而当销售公司已知会售罄时，则会出现垄断现象。

基线状态下不存在分离均衡（命题2）。以下关于汇集和半分离的特性，是同一PBE模型的替代结果，并非对一个策略组合的同时限制。保留了对称汇集版本和原有的全有或全无平局规则。

对于对称池化均衡（命题3，物理部分~13–14），
\[
p_{i1}^{c}=p_{i1}^{u}=p_1^{*}=\alpha,
\qquad i=1,2,
\qquad \theta(\alpha)=\alpha.
\]
将(1)--(2)代入后得到
\[
F^{S}(p_2)=1-\frac{\alpha(1-p_2)}{(1-\alpha)p_2},
\qquad
F^{N}(p_2)=1-\frac{\alpha}{p_2},
\qquad p_2\in[\alpha,1),
\]
在$F^{N}$中的$1$位置有一个大小为$\alpha$的原子。这些表达式继承了内部-后部前提$\theta=\alpha\in(0,1)$。在给定类型的情况下，预期的总利润为
\[
\text{constrained firm's profit}=\alpha,
\qquad
\text{unconstrained firm's profit}=\frac{3\alpha}{2}.
\]
每家公司的预期事前利润是
\[
\alpha\left[\alpha+\frac{3}{2}(1-\alpha)\right]
 =\alpha\frac{3-\alpha}{2}.
\]
对于在$p_1<\alpha$下进行无约束的低价竞争的情况，让$\theta'$成为由此产生的非均衡信念。论文中给出了…
\[
\text{deviation payoff}=p_1+\theta',
\qquad
\theta'\leq\frac{\alpha}{2}
\]
将其视为一种有效的约束条件，以阻止此类偏差的发生。不过，这并非一种必然的约束条件。在对手的报价路径上进行的贝叶斯更新仍然会得出$\alpha$的结果，即使出现向上的偏差。其他历史数据则仍然遵循上述共同的顺序最优性和信念要求。

对于半分离均衡的情况（物理部分：~14–16），$\theta(p)$表示以第一期价格$p$收费的企业的后验概率，而$Q(p)$则表示该企业在第1期销售产品的概率。正文部分则直接引用了与类型相关的预期利润。
\begin{equation}
\Pi^{c}(p)=Q(p)p+
\int_{\underline p}^{p}\theta(p')\,d\bigl(1-Q(p')\bigr),
\tag{3}
\end{equation}
\begin{equation}
\Pi^{u}(p)=Q(p)\bigl(p+\theta(p)\bigr)+
\int_{\underline p}^{p}\theta(p')\,d\bigl(1-Q(p')\bigr).
\tag{4}
\end{equation}
紧接着第(3)段之后的正文段落中，将$\underline p$称为\emph{第二个}周期内的最低价格。这一描述被原样保留了下来。而附录A则使用了\emph{第一个}周期的支持区间的下限，并将其确定为$\alpha$。这两种描述在这里并未被明确区分。下面的具体描述则采用了附录中给出的下限为$\alpha$的支付行。

对于共同支持中的所有价格而言，每种类型的均衡利润必须保持恒定。相关的收益关系如下：
\[
\Pi^{u}(p)=\Pi^{c}(p)+Q(p)\theta(p).
\]
因此，$Q(p)\theta(p)$在支持范围内是恒定的且为正值；后验概率会随着价格的增加而上升。一个不受约束的 firm 只会使用那些在第一阶段有正销售概率的价格。上端点有一个质量点。

命题 4 通过以下方式对半分离家族进行索引：
\[
\bar p\in(\alpha,1].
\]
这两种类型都会随机化它们在第一阶段的价格。
\[
[\alpha,\bar p),
\]
在 $\bar p$ 位置有一个质量点，而第二周期的价格范围则如下所述：
\[
[\theta(\alpha),1].
\]
这里的$\bar p$是一个均衡指数，而不是额外的外部价格-行为约束。具有约束条件的类型的均衡利润是$\alpha$。

附录中选择的内部结构（物理页面~38–40）适用于每一个$p\in[\alpha,\bar p)$。其中，$\tilde p$表示对手方第一阶段的报价。
\[
Q(p)=\Pr(\tilde p>p),\qquad p\in[\alpha,\bar p),
\]
这与正文中的销售概率相符，且远离了上面的原子。
\[
\Pi^{c}(p)=Q(p)p+
\int_{\alpha}^{p}-Q'(\tilde p)\theta(\tilde p)\,d\tilde p,
\qquad p\in[\alpha,\bar p),
\]
\[
\begin{aligned}
\Pi^{u}(p)
 &=Q(p)\bigl(p+\theta(p)\bigr)
   +\int_{\alpha}^{p}-Q'(\tilde p)\theta(\tilde p)\,d\tilde p\\
 &=\Pi^{c}(p)+Q(p)\theta(p),
\qquad p\in[\alpha,\bar p).
\end{aligned}
\]
定义性的冷漠和边界关系包括：
\[
\Pi^{c}(p)=\alpha,
\qquad
\Pi^{u}(p)=\Pi^{u}(\alpha),
\qquad p\in[\alpha,\bar p),
\]
\[
Q(\alpha)=1,
\qquad
Q(p)\theta(p)=\theta(\alpha),
\qquad p\in[\alpha,\bar p),
\]
\[
\theta(\bar p)=\bar p,
\qquad
\lim_{p\uparrow\bar p}\theta(p)
 =\frac{1}{2}\theta(\bar p)=\frac{\bar p}{2}.
\]
最终得到的后验值为：
\[
\theta(p)=
\frac{p}{-W_{-1}\!\left(-\dfrac{2p}{e^{2}\bar p}\right)},
\qquad p\in[\alpha,\bar p),
\]
并且，特别是
\[
\theta(\alpha)=
\frac{-\alpha}{W_{-1}\!\left(-\dfrac{2\alpha}{e^{2}\bar p}\right)}.
\]
源代码中将$W_{-1}$称为Lambert-$W$函数的负分支，并引用了相关恒等式。
\[
W(xe^{x})=x.
\]
附录中列出了那些偏离路径的信念。
\[
\theta(p)=0\quad\text{for }p<\alpha,
\qquad
\theta(p)=1\quad\text{for }p>\bar p.
\]
对于$[\alpha,\bar p]$中的价格，信念是通过贝叶斯更新从均衡策略中获得的。总不受限制的利润则通过来源的身份来保留。
\[
\Pi^{u}(\alpha)=\alpha+Q(\alpha)\theta(\alpha),
\qquad Q(\alpha)=1,
\]
与上述印刷后的公式一起。第一时期类型依赖的混合策略仍然遵循贝叶斯规则以及常见的PBE最优性条件；没有引入额外的类型特定密度公式。

\emph{打印利润冲突。}关于物理p的命题4。~15种无约束类型的利润打印结果如下：
\[
\alpha\left(
1-\left[
W_{-1}\!\left(\frac{2\alpha}{e^{2}\bar p}\right)
\right]^{-1}
\right).
\]
正参数在提供的原始页面图像中可见。对于正$\alpha$和$\bar p$，它与所述的负实数分支以及同一页面后文中的负参数不兼容。附录A中，物理p.~40单独列出。
\[
\Pi^{u}(\alpha)
 =\alpha+Q(\alpha)\theta(\alpha)
 =\alpha+\frac{\alpha}{W\!\left(-\dfrac{2\alpha}{e^{2}\bar p}\right)}.
\]
其最后的表达式也与前面的恒等式以及附录中的负分支后项相矛盾。这两种字面意义上的闭式形式被记录为打印冲突，而不是作为额外的支付等式来施加。无论是符号还是分支都没有被自动修复。

\emph{印刷的上原子表示法。}附录A在质量点计算中使用了$Q(\bar p)$，具体方法如下：
\[
\begin{aligned}
&\lim_{p\uparrow\bar p}Q(p)p
 +\int_{\alpha}^{\bar p}-Q'(\tilde p)\theta(\tilde p)\,d\tilde p\\
&\qquad=
 \frac12\bigl(Q(\bar p)\bar p+Q(\bar p)\theta(\bar p)\bigr)
 +\int_{\alpha}^{\bar p}-Q'(\tilde p)\theta(\tilde p)\,d\tilde p,
\end{aligned}
\]
\[
\lim_{p\uparrow\bar p}Q(p)\theta(p)
 =\frac12Q(\bar p)\theta(\bar p),
\qquad
Q(\bar p)=\frac{2\theta(\alpha)}{\bar p}.
\]
这些是附录中的端点计算。其中的$Q(\bar p)$作为完整的竞争价格原子：然后，$1/2$因素会应用平局规则。这与主文本中的销售概率$\bar p$不同，后者包括赢得平局的情况，也与严格尾部表达式$\Pr(\tilde p>\bar p)$不同，后者在支持端点处为零。这些不兼容的端点含义不会同时应用于主文本的$Q$。上述内部定义和后期边界关系被保留，不会用极限值替换端点值。

在附录中的贝叶斯归一化计算中，也出现了相同的端点用法（物理 p.~41）：
\begin{equation}
\int_{\alpha}^{\bar p}-Q'(p)\theta(p)\,dp
 +Q(\bar p)\bar p
 =-\theta(\alpha)
 W_{-1}\!\left(-\frac{2\alpha}{e^{2}\bar p}\right)
 =\alpha.
\tag{11}
\end{equation}
这是打印出来的预期后验计算结果，其中附录中使用了$Q(\bar p)$来表示完整原子的情况；它并没有在原子级别定义主要文本的销售概率。相反，前面的物理页面则将其作为提出的先验一致性条件来打印出来。
\[
\int_{\alpha}^{\bar p}-Q'(p)p\,dp
 +Q(\bar p)\bar p=\alpha.
\]
这里的被积函数中包含$p$，而(11)则包含$\theta(p)$。这些都是不同的打印表达式。后一页中的表达式并不强制遵循贝叶斯一致性原则；所选择的链遵循后验加权表达式(11)，同时明确保留了端点表示法的限制。

\emph{允许端点。}源包括$\bar p=1$，并需要$\theta(\bar p)=\bar p$。如果其声明两种类型都有质量点，而被理解为每种类型都需要严格正的原子，那么当两种类型的先验概率均为正时，就会在$\bar p=1$处与贝叶斯更新产生冲突。打印出的参数范围和原子声明保持不变；没有添加替代原子规范或限制$\bar p<1$。同样，公式(1)--(2)在此端点仍被限制为$0<\theta<1$。

\medskip
\noindent\emph{来源范围。}本文采用了所提供的主论文设置和PBE定义、基线延续性及均衡特性描述，以及附录A中的相关定义段落。物理PDF页面与论文中的印刷页码是分开的。比较静态证明、完整信息基准测试以及扩展内容并非本公式的附加条件。引用的在线补充材料和附录C-D并未提供，因此未被验证；其中没有任何行被使用。

\end{document}
